\documentclass[10pt,a4paper]{amsart}
\usepackage[margin=19mm]{geometry}
\usepackage{amsmath,amssymb,mathtools}
\usepackage[scr=rsfs,cal=euler]{mathalfa}
\usepackage{xfrac}
\usepackage[unicode,hidelinks,bookmarks=false]{hyperref}
\newcommand{\ie}{i.e.\ }
\newcommand{\cf}{cf.\ }
\newcommand{\resp}{respectively\ }
\newcommand{\Subsection}{\subsection}
\providecommand{\nobreakdash}{\nobreak\hbox{-}}
\setlength{\emergencystretch}{2em}
\multlinegap0pt
\usepackage{bm}
\usepackage{bbm}
\usepackage{dsfont}
\usepackage{xspace}
\usepackage{cancel}
\usepackage{mathtools}
\usepackage{amsthm}
\makeatletter
\@ifundefined{theorem}{\newtheorem{theorem}{Theorem}}{}
\@ifundefined{lemma}{\newtheorem{lemma}[theorem]{Lemma}}{}
\makeatother
\title[Finite groups with nearly half as many cyclic subgroups as e]{Finite groups with nearly half as many cyclic subgroups as elements}
\author{Vaibhav Chhajer, Sumana Hatui, Palash Sharma}
\date{Source: arXiv:2506.21163v2 (CC BY 4.0)}
\begin{document}
\maketitle

\noindent\emph{Source and licence.} Finite groups with nearly half as many cyclic subgroups as elements. Version arXiv:2506.21163v2 (CC BY 4.0), distributed under the Creative Commons Attribution 4.0 International licence, as recorded in the arXiv licence metadata for this exact version and shown on the abstract page https://arxiv.org/abs/2506.21163v2. Full rights notice: licenses/arxiv-2506.21163-CC-BY-4.0.txt.\par\smallskip

\noindent\textit{Editorial scope.} Counted unit: the Lemma whose source label is \texttt{lemma5} in the article (source0.tex lines 464--477, source label \texttt{lemma5}), delivered together with the article's own complete original proof of it (source0.tex lines 478--496, one \texttt{proof} environment of 19 lines). No mathematics is written, repaired or re-derived: statement and proof are the article's own text line for line. Required context retained uncounted: none. Every definition, display and label of those lines is retained unchanged. Omitted from this package and not redistributed: 1--463, 497--846. Nothing inside a retained excerpt is omitted or altered: every retained body is delivered exactly as the article's lines are written, so no omission receipt of any kind accompanies this package. Citations: the retained excerpts carry no citation command at all, so no bibliography entry of the article is transcribed and no reference is redistributed. Reference adaptation: none was needed and none was made; every reference inside the delivered unit resolves inside it, so no unresolved reference or citation is produced. Printed slips kept literal: no printed word, symbol, hypothesis or number of the counted statement or of its proof was altered. Layout and class: the article's own class is \texttt{amsart} with packages including amssymb, amscd, amsmath, eucal, array, float, xy, pdflscape, hyperref, enumerate, tikz, tikz-cd, relsize, geometry; the wrapper is the lane's pinned \texttt{amsart} editorial wrapper at 10pt on A4, which restates the theorem-like environments the retained text needs and adds the article's own macro definitions that the retained text needs (none), with extra wrapper packages (bm, bbm, dsfont, xspace, cancel, mathtools). The delivered numbering is the unit's own. Assets: no retained span contains a figure environment, a graphics-inclusion command, a PDF-page inclusion, a table or a TikZ picture; no image file is copied into this package, the package declares no asset dependency and no image file is redistributed with it. Licence: version arXiv:2506.21163v2 of this article is distributed under the Creative Commons Attribution 4.0 International licence, as recorded in the arXiv licence metadata for this exact version and shown on the abstract page https://arxiv.org/abs/2506.21163v2. A full rights notice is delivered at licenses/arxiv-2506.21163-CC-BY-4.0.txt. Characters: every retained excerpt is pure ASCII, so the delivered engine, XeLaTeX with its default Latin Modern text font, reports no missing character.
\begin{lemma}\label{lemma5}
   There is no group $G$ satisfying $(P_2)$ and one of the following conditions.
   \begin{enumerate}
\item $K_1$ or $K_2 \cong \mathbb Z_{10}$.

\item $K_1$ or $K_2 \cong \mathbb Z_{5}$.



        \item $K_1 \cong \mathbb{Z}_{8}$ and $K_2 \cong \mathbb{Z}_{12}$.


   \end{enumerate}
\end{lemma}

\begin{proof} 
$(1)$ If $K_1\cong \mathbb{Z}_{10}$, then $K_2\cong \mathbb{Z}_5$ such that $K_2 < K_1$. Taking an element, say $x_2$, of order $2$ such that $x_2\notin K_1$, we get a subgroup $K_2\langle x_2\rangle $ of order $10$, contradicting $(P_2)$.

    
    $(2)$ Suppose $K_1\cong \mathbb{Z}_5$. If $K_2\ncong \mathbb{Z}_5$, then $K_1\trianglelefteq G$. Thus, the group will contain either $D_{10}$ or $\mathbb{Z}_{10}$, a contradiction.
    

    Now assume, $K_2\cong \mathbb{Z}_5$. Then, using the previous argument, we conclude that $K_i$ are not normal in $G$, and they are conjugates of each other. 
    If any element of order $2$ lies inside $N(K_1)$, then as $K_1\trianglelefteq N(K_1)$, we will get a subgroup of order $10$,  contradicting $(P_2)$. Also note that $3 \nmid |N(K_1)|$ as we will have an element of order $15$, contradicting $(P_{2})$. Thus $N(K_1)$ must be a $5$-group. Therefore, $|N(K_1)| = 5^k$. 
    Now, if $k>1$, we will get a subgroup of order $25$ which would be either $\mathbb{Z}_5\times \mathbb{Z}_5$ or $\mathbb{Z}_{25}$, both are not possible. Hence, $|N(K_1)| = 5 \implies |G| = 2|N(K_1)| = 10$, and no such group of order $10$ satisfies $|C(G)| = \frac{|G|}{2}$.

% Suppose $K_1\cong \mathbb{Z}_5$. If $K_2\ncong \mathbb{Z}_5$ then the group will contain either $D_{10}$ or $\mathbb{Z}_{10}$ as $K_1\trianglelefteq G$ contradicting $(P_2)$. Hence $K_2\cong \mathbb{Z}_5$,  now in this case, if any one of $K_1$ or $K_2$ is normal in $G$, then previous argument will follow. So assume that both are not normal in $G$, which means that these two subgroups are conjugates of each other and moreover they completely exhaust the list of conjugates of $K_1$, thus index of normalizer of $K_1$ is then $|G:N(K_1)|=2$. If any element of order $2$ lies inside $N(K_1)$ then since $K_1\trianglelefteq N(K_1)$ we will get a subgroup of order $10$ which would be either $D_{10}$ or $Z_{10}$ contradicting $(P_2)$. Also notice that $3 \nmid |N(K_1)|$, because then we will have an element of order $15$, contradicting $(P_{2})$. Thus $N(K_1)$ must be a $5$-group. Therefore $|N(K_1)| = 5^k$. Now if $k>1$, we will get a subgroup of order $25$ which would be either $\mathbb{Z}_5\times \mathbb{Z}_5$ or $\mathbb{Z}_{25}$, the former is not possible as it contains more than two subgroups isomorphic to $\mathbb{Z}_5$ and the later is not possible as it contains an element of order $25$ contradicting $(P_2)$. Hence, $|N(K_1)| = 5 \implies |G| = 2|N(K_1)| = 10$ and no such group  $10$ which satisfies $C(G) = \frac{|G|}{2}$.
    




$(3)$ In this case, for some $i\in\{1,\cdots,n-6\}$,  we have $H_i< K_2\leq G$  such that $H_{i} \cong \mathbb{ Z}_3$. Clearly $K_1\trianglelefteq G$ and $K_1 \cap H_{i}= \{e\}$, so $K_1H_i\leq G$ with $K_1H_i\cong \mathbb{Z}_{8}\rtimes \mathbb{Z}_{3}\cong \mathbb{Z}_{24}$, a contradiction to $(P_2)$.
\end{proof} 

\end{document}
